\chapter*{Lời Tựa}
\addcontentsline{toc}{chapter}{Lời Tựa}
\markboth{Lời Tựa}{Lời Tựa}

\begin{truyen}{Lời Tựa}{A book of educational situations for teachers, students, and parents}
\chuhoa{T}{rong} mỗi ngôi trường, người ta thường gọi tên rất nhanh: học sinh ngoan hay học sinh hư, giáo viên tốt hay giáo viên xấu.
\textit{(In every school, labels come quickly: a good student or a bad one, a good teacher or a harmful one.)}

Nhưng thực tế phức tạp hơn nhiều. Có người thầy nghiêm mà cứu được cả một chặng đời của học sinh; cũng có người nói nhiều về tử tế nhưng lại để lại những vết thương khó gọi tên. Có học sinh từng quậy phá, quay cóp, vô lễ, nhưng rồi lớn lên thành người biết tự trọng và có ích.
\textit{(Reality is more complex. A strict teacher may still rescue a student's future; a troubled student may still grow into a decent, responsible adult.)}

Cuốn sách này không viết để kết án. Nó được hình dung như một tấm gương nhiều mặt, nơi cả thầy lẫn trò đều có thể soi lại cách mình dạy, cách mình học, cách mình nói, và cách mình sửa sai.
\textit{(This book is not written to condemn. It is written as a many-sided mirror for both teachers and students to examine how they teach, learn, speak, and recover from mistakes.)}
\end{truyen}

\begin{baihoc}
Điểm tựa lớn nhất của giáo dục không nằm ở việc phân loại ai xấu, ai tốt, mà ở khả năng mở ra một con đường sửa mình cho cả hai phía.
\textit{(The deepest strength of education is not classification, but the ability to open a path of change.)}
\end{baihoc}

\begin{ghinhoanh}
Labels can hide the truth.

Education should correct without crushing.

Every situation can become a turning point.
\end{ghinhoanh}

\begin{tuhocnhanh}
1. Trong lớp học, điều gì dễ làm học sinh xấu hổ nhất? \textit{(What most easily shames a student in class?)}

2. Một người thầy tốt có nhất thiết phải luôn mềm mỏng không? \textit{(Must a good teacher always be gentle?)}

3. Một học sinh từng sai có cần bị nhìn bằng quá khứ mãi không? \textit{(Should a student always be judged by the past?)}
\end{tuhocnhanh}

\begin{center}
\textit{\color{truyengold}``Mỗi tình huống nhỏ trong lớp học đều có thể thành một bài học lớn của cuộc đời.''}
\end{center}

\ngancach
